\chapter{Plano de trabalho}
\label{cap:plano}

% wiki: Roadmap; diario/Semana 1..4; log.md (M1, M2)

\section{Organização em frentes}
\label{sec:divisao}

O trabalho está dividido em três frentes, que avançam em paralelo graças aos
contratos da Seção~\ref{sec:contratos}:

\begin{itemize}
  \item Frente A, análise estática (Pedro Henrique Machado):
    compilação, estágios 1 a 3, registro de contexto e resolvedor de pedidos de
    contexto.
  \item Frente B, integração com \llm (Lucas de Oliveira Martim):
    triagem, desenho de \textit{prompt}, ablação, reparo e sua avaliação.
  \item Frente C, painel: interface de operação, construída no tempo
    livre das outras duas e ordenada pelo retorno imediato de cada parte.
\end{itemize}

A frente B não espera pela frente A. Até a integração, ela trabalha sobre os 20
registros de contexto montados à mão (Seção~\ref{sec:benchmarks}) e sobre um
resolvedor simulado, que é trocado pelo resolvedor real sem alterar mais nada.

\section{Cronograma}
\label{sec:cronograma}

O plano cobre dez semanas, de 31 de agosto a 6 de novembro de 2026
(Tabela~\ref{tab:cronograma}).

\begin{table}[htb]
  \centering
  \caption{Cronograma das frentes A e B}
  \label{tab:cronograma}
  \small
  \begin{tabular}{@{}lL{6.2cm}L{6.2cm}@{}}
    \toprule
    Semana & Frente A (análise estática) & Frente B (\llm) \\
    \midrule
    1 & contratos; compilação dos 31 casos & contratos; 20 registros; sonda de viabilidade \\
    2 & gabarito certificado; unidade e regra de casamento & ambiente de execução e \textit{prompt} simples \\
    3--4 & estágio 1: candidatos & regras do domínio; pedidos de contexto \\
    5 & estágio 2: prioridade e mascaramento & decomposição e autovalidação; ablação \\
    6--7 & estágio 3: viabilidade com Z3 & votação, consistência e reparo \\
    8 & integração & integração com o resolvedor real \\
    9 & suíte de programas reais & triagem e reparo na suíte real \\
    10 & consolidação e escrita & consolidação e escrita \\
    \bottomrule
  \end{tabular}
  \fonte{Elaborada pelos autores.}
\end{table}

Cinco marcos organizam o acompanhamento (Tabela~\ref{tab:marcos}).

\begin{table}[htb]
  \centering
  \caption{Marcos do projeto}
  \label{tab:marcos}
  \small
  \begin{tabular}{@{}llL{10.5cm}@{}}
    \toprule
    Marco & Data & Entrega \\
    \midrule
    M1 & 11/09 & contratos congelados; gabarito certificado; 20 registros; resposta da sonda de viabilidade do \llm \\
    M2 & 25/09 & estágio 1 nos 31 casos do Racebench, com portão de recall 48/48 \\
    M3 & 02/10 & tabela de ablação da triagem \\
    M4 & 23/10 & ciclo completo no Racebench: análise, Z3, triagem e reparo \\
    M5 & 06/11 & suíte de programas reais; métricas finais; texto completo \\
    \bottomrule
  \end{tabular}
  \fonte{Elaborada pelos autores.}
\end{table}
